% sec-01-introduction.tex - Section 1, Introduction.  FINAL stage.
% Figures 1 and 2 drawn; Tables 4 and 5 populated.  Senior-author pass: the
% Action column of Table 4 gains 0.10 cm from the Date column, which needs only
% eleven characters, and a closing subsection states what the paper claims and
% what it does not.
% Sources: new-trial-system/references/trump-ai-cancer-2025-2026.bib for every
% citation key, and new-trial-system/d2/fig-02 for every cell of Table 5.
\section{Introduction}
\label{sec:introduction}

Between January 2025 and July 2026 the executive branch committed to artificial
intelligence and to cancer at the same time, and often in the same document.
Executive Order 14179 removed barriers to American leadership in AI
\cite{eo14179}, and Executive Order 14212 established the MAHA Commission
\cite{eo14212}. Executive Order 14355, signed September 30, 2025, is the only
executive order aimed squarely at both topics: it tasks the MAHA Commission, in
coordination with HHS, the Assistant to the President for Science and
Technology, and the Special Advisor for AI and Crypto, with developing ways to
use AI for improved diagnoses, treatments, cures, and prevention of pediatric
cancer, beginning with the Childhood Cancer Data Initiative and with predictive
modeling of patient response, disease progression, and treatment toxicity
\cite{eo14355, eo14355app, akin_eo14355}. The same day, HHS raised the
Initiative's budget from \$50 million to \$100 million and committed to bringing
private-sector partners in to apply advanced AI \cite{ccdidoubling_nih}. The
MAHA Strategy of September 9, 2025 had already directed NIH and OSTP to develop
an AI-driven approach to transform pediatric cancer research and clinical trials
\cite{maha_strategy}.

The machinery followed the mandate. Executive Order 14363 launched the Genesis
Mission on November 24, 2025 \cite{eo14363, doe_genesis_launch}, and its NIH arm,
the Bio Genesis Mission announced July 22, 2026, carries more than \$1.2 billion
in obligated FY2026 and planned FY2027 funding across predicting the behavior of
living systems, accelerating drug translation, and unlocking new pediatric
cancer treatments \cite{biogenesis_nih, hpcwire_genesis}. Inside the agencies,
the FDA fielded Elsa on June 2, 2025, a generative tool to help staff synthesize
information for tasks including clinical protocol review, and followed it with
an agency-wide agentic platform on December 1, 2025
\cite{crs_ai_healthcare, hhs_ai_strategy}. CMS announced the WISeR model and,
at the July 30, 2025 health technology event, an interoperability framework that
was later institutionalized in a dedicated technology office
\cite{cms_healthtech, cms_office_healthtech}. On June 22, 2026 HHS launched a
department-wide effort to restore American leadership in clinical trials, of
which Operation TrialBlazer, NCI working with cancer centers to streamline trial
activation and improve enrollment, is the oncology arm
\cite{trialblazer, hhstrialblazer2026}. Table 4 places ten of these actions on
one chronology with what each supplies to a sponsor.

\begin{apptable}
\begin{tabularx}{\textwidth}{>{\raggedright\arraybackslash}p{5.35cm} >{\raggedright\arraybackslash}p{1.95cm} >{\raggedright\arraybackslash}X}
\headrow \thead{Action} & \thead{Date} & \thead{What it supplies to a pancreatic cancer sponsor}\\
\midrule
EO 14179, Removing Barriers to American Leadership in AI & Jan 23, 2025 & Policy cover for AI inside a regulated submission \\
EO 14212, MAHA Commission & Feb 13, 2025 & The parent authority EO 14355 cites \\
FDA Elsa generative AI tool & Jun 2, 2025 & An agency that already reads AI-assisted protocol text \\
MAHA Strategy & Sep 9, 2025 & A direction to transform trials with AI \\
EO 14355, Unlocking Cures for Pediatric Cancer with AI & Sep 30, 2025 & The only order aimed at AI and cancer together \\
CCDI budget doubled to \$100 million & Sep 30, 2025 & Money for AI-driven oncology data work \\
EO 14363, Launching the Genesis Mission & Nov 24, 2025 & An integrated Federal AI platform \\
HHS AI Strategy and the agentic platform & Dec 1 to 4, 2025 & Agency-side capacity to process AI-assisted filings \\
HHS clinical trial effort and Operation TrialBlazer & Jun 22, 2026 & Faster trial activation at cancer centers \\
Bio Genesis Mission & Jul 22, 2026 & Over \$1.2 billion toward AI in biology \\
\end{tabularx}
\end{apptable}
\vspace{-0.6cm}
\tabcap{Table 4. Ten Federal actions between January 2025 and July 2026, and what\\each
supplies to a sponsor preparing a pancreatic cancer trial dossier.}

The commitment is real and it is funded, and it ran alongside proposed
reductions: the FY2027 request of April 1, 2026 proposed roughly \$41.1 billion
for NIH against \$48.7 billion in FY2026, though NCI would receive \$7.353
billion, an increase of \$8.9 million, and Congress added \$128 million for
cancer research in the final FY2026 bill
\cite{aacr_policymonitor_2026, cancerdiscovery_fy27, science_fundingbill}. What
none of it supplies is a sponsor-side method. Compute, data, agency review
capacity, and faster site activation are all being built. The document a sponsor
must file is still written the way it was written in 2015, and that is the gap
Figure 1 states and this paper closes.

% ---------------------------------------------------------------------------
% FIGURE 1.  mermaid-type, flowchart LR.  Eleven Federal actions in three
% clusters, three capabilities, one gap, one system.  Absolute coordinates
% throughout.  The only two edges that could cross a cluster, A3 and A2 to C3,
% are routed at y = -2.95 below every frame.
% ---------------------------------------------------------------------------
\begin{appfloat}[!tb]
\begin{appfig}[mermaid-type / flowchart LR]
% --- cluster 1, authority
\node[mmsoft,text width=27mm] (a1) at (0,3.15) {EO 14179\\remove barriers to\\American AI leadership};
\node[mmsoft,text width=27mm] (a2) at (0,1.65) {EO 14212\\MAHA Commission,\\the parent authority};
\node[mmmid,text width=27mm]  (a3) at (0,0.15) {EO 14355\\AI for pediatric cancer,\\the one order at both};
\node[mmsoft,text width=27mm] (a4) at (0,-1.35) {EO 14363\\Genesis Mission,\\integrated AI platform};
\begin{scope}[on background layer]
\node[mmlane,fit=(a1)(a2)(a3)(a4),inner sep=6pt] (cl1) {};
\end{scope}
\node[mmlanetitle,anchor=south west] at ([yshift=1.3mm]cl1.north west) {Authority, Jan 2025 to Dec 2025};
% --- cluster 2, money and machinery
\node[mmsoft,text width=27mm] (m1) at (4.10,3.15) {CCDI doubled\\50M to 100M dollars};
\node[mmmid,text width=27mm]  (m2) at (4.10,1.65) {Bio Genesis Mission\\above 1.2B obligated};
\node[mmsoft,text width=27mm] (m3) at (4.10,0.15) {FDA Elsa, then an\\agency wide agentic\\AI platform};
\node[mmsoft,text width=27mm] (m4) at (4.10,-1.35) {HHS AI Strategy\\and CMS WISeR};
\begin{scope}[on background layer]
\node[mmlane,fit=(m1)(m2)(m3)(m4),inner sep=6pt] (cl2) {};
\end{scope}
\node[mmlanetitle,anchor=south west] at ([yshift=1.3mm]cl2.north west) {Money and machinery, 2025 to 2026};
% --- cluster 3, trial machinery
\node[mmmid,text width=27mm]  (t1) at (8.20,2.40) {Operation TrialBlazer\\faster trial activation};
\node[mmmid,text width=27mm]  (t2) at (8.20,0.90) {HHS department wide\\clinical trial effort};
\node[mmsoft,text width=27mm] (t3) at (8.20,-0.60) {MAHA Strategy\\AI driven pediatric\\trial transformation};
\begin{scope}[on background layer]
\node[mmlane,fit=(t1)(t2)(t3),inner sep=6pt] (cl3) {};
\end{scope}
\node[mmlanetitle,anchor=south west] at ([yshift=1.3mm]cl3.north west) {Trial machinery, 2026};
% --- capabilities
\node[mmin,text width=25mm] (c1) at (11.85,2.40) {Supplied\\compute, data,\\and platform};
\node[mmin,text width=25mm] (c2) at (11.85,0.55) {Supplied\\review capacity\\inside the agency};
\node[mmin,text width=25mm] (c3) at (11.85,-1.30) {Supplied\\faster activation\\at cancer centers};
% --- gap and system
\node[mmgrayd,text width=33mm,minimum height=15mm] (gap) at (15.55,0.55)
  {Not supplied\\a sponsor side method that\\produces an IND, a protocol,\\a bill and a funding file\\on a 1 to 4 day scale};
\node[mmgoal,text width=33mm] (sys) at (15.55,-2.55) {This paper\\the new PDAC trial system};
% --- edges, cluster to capability
\draw[mmedge] (a1.east) -- (c1.west);
\draw[mmedge] (a4.east) to[out=0,in=200] (c1.west);
\draw[mmedge] (m1.east) -- (c1.west);
\draw[mmedge] (m2.east) -- (c1.west);
\draw[mmedge] (m3.east) -- (c2.west);
\draw[mmedge] (m4.east) to[out=0,in=195] (c2.west);
\draw[mmedge] (t1.east) -- (c3.west);
\draw[mmedge] (t2.east) -- (c3.west);
\draw[mmedge] (t3.east) -- (c3.west);
% the two edges that must clear the money cluster: routed below every frame
\draw[mmedge] (a3.south) -- ++(0,-1.55) -- (2.05,-2.95) -- (11.85,-2.95) -- (c3.south);
\draw[mmedge] (a2.west) -- ++(-0.45,0) -- (-2.05,-2.95) -- (11.85,-2.95);
% --- edges, capability to gap
\draw[mmedged] (c1.east) -- (gap.west);
\draw[mmedged] (c2.east) -- (gap.west);
\draw[mmedged] (c3.east) -- (gap.west);
\draw[mmedgeb,line width=1.1pt] (gap.south) -- (sys.north);
\pnote{-2.05}{-4.05}{Eleven actions, three capabilities, one gap. Every action is cited to
\texttt{trump-ai-cancer-2025-2026.bib}; no action is\\counted twice, and the gap
node is the only Mist Gray fill on the canvas because it is the figure's
subject.}
\end{appfig}
\vspace{-0.6cm}
\figcaption{Figure 1. Eleven Federal AI and cancer actions from January 2025 to
July\\2026, the three capabilities they supply, and the one they leave unfilled.}
\end{appfloat}

Table 5 and Figure 2 carry the same ten axes: the first as text a reader can
cite, the second as a grid a reader can scan. The prior system is not described
here as bad work. It is described as work whose unit is the team-month, and a
team-month cannot produce an IND, two protocols, five bill versions and fourteen
funding artifacts inside a single quarter. The remainder of this paper is the
evidence for that sentence: \S\ref{sec:methods} gives the method, \S\ref{sec:ind} through \S\ref{sec:peerreview} give the four
deliverables and the review regime that cleared them, and \S\ref{sec:limitations} states what
remains exposed.

% ---------------------------------------------------------------------------
% FIGURE 2.  d2-type, grid.  A true grid: one row height, three fixed column
% widths, no edges, because the comparison is positional.
% ---------------------------------------------------------------------------
\begin{appfloat}[!tb]
\begin{appfig}[d2-type / grid]
\def\rh{0.62}
% header
\node[d2cellh,minimum width=4.6cm,minimum height=\rh cm,anchor=north west] at (0,0) {Operating axis};
\node[d2cellh,minimum width=5.0cm,minimum height=\rh cm,anchor=north west] at (4.6,0) {Prior system};
\node[d2cellh,minimum width=5.0cm,minimum height=\rh cm,anchor=north west] at (9.6,0) {New system};
\foreach \i/\ax/\op/\np in {%
  1/{Document production unit}/{Team-month}/{Prompt-hour},
  2/{IND assembly time}/{Months across a group}/{Four days, one author},
  3/{Peer review entry point}/{After completion}/{During development},
  4/{Review latency}/{7 to 8 weeks best case}/{Same day, hour scale},
  5/{Reviewers per round}/{Two to three humans}/{Three model manufacturers},
  6/{Virtual trial cost}/{Above 120000 dollars}/{36330 dollars projected},
  7/{Provenance record}/{Acknowledgement paragraph}/{Commit history plus DOI},
  8/{Figure format}/{Raster, not reusable}/{Machine readable, reusable},
  9/{Revision granularity}/{Whole manuscript}/{One file, one commit}}{
  \node[d2cellg,minimum width=4.6cm,minimum height=\rh cm,anchor=north west]
    at (0,{-\rh*\i}) {\ax};
  \node[d2cell,minimum width=5.0cm,minimum height=\rh cm,anchor=north west]
    at (4.6,{-\rh*\i}) {\op};
  \node[d2celll,minimum width=5.0cm,minimum height=\rh cm,anchor=north west]
    at (9.6,{-\rh*\i}) {\np};}
% row 10, the emphasis row
\node[d2cellg,minimum width=4.6cm,minimum height=\rh cm,anchor=north west] at (0,{-\rh*10}) {Author count};
\node[d2cell,minimum width=5.0cm,minimum height=\rh cm,anchor=north west] at (4.6,{-\rh*10}) {Group, plus a CRO};
\node[d2cellk,minimum width=5.0cm,minimum height=\rh cm,anchor=north west] at (9.6,{-\rh*10}) {One, plus the models};
% header rule
\draw[trialchar,line width=0.7pt] (0,{-\rh}) -- (14.6,{-\rh});
\node[d2mid,text width=44mm,anchor=north] at (7.30,{-\rh*10-0.85})
  {Row 10 is the axis the other nine follow from};
\pnote{0}{-7.75}{Every cell carries a measured value from a deposited source, not an adjective.
The prior column is white and the new\\column pale burgundy; the one emphasis
cell is the axis that produces the other nine.}
\end{appfig}
\vspace{-0.6cm}
\figcaption{Figure 2. Ten operating axes on which the prior and the new trial system\\
differ, each cell carrying a measured value from the deposited record.}
\end{appfloat}

\begin{apptable}
\begin{tabularx}{\textwidth}{>{\raggedright\arraybackslash}p{3.35cm} >{\raggedright\arraybackslash}p{3.85cm} >{\raggedright\arraybackslash}p{3.85cm} >{\raggedright\arraybackslash}X}
\headrow \thead{Operating axis} & \thead{Prior system} & \thead{New system} & \thead{Established in}\\
\midrule
Document production unit & Team-month & Prompt-hour & \S\ref{sec:methods} \\
IND assembly time & Months across a group & Four days, one author & \S\ref{sec:ind} \\
Peer review entry point & After completion & During development & \S\ref{sec:peerreview} \\
Review latency & 7 to 8 weeks best case & Same day & \S\ref{sec:peerreview} \\
Reviewers per round & 2 to 3 humans & 3 model manufacturers & \S\ref{sec:peerreview} \\
Virtual trial cost & Above \$120,000 per run & \$36,330 projected & \S\ref{sec:funding} \\
Provenance record & Acknowledgement paragraph & Commit history plus DOI & \S\ref{sec:methods} \\
Figure format & Raster, not reusable & Machine readable, reusable & \S\ref{sec:methods} \\
Revision granularity & Whole manuscript & One file, one commit & \S\ref{sec:methods} \\
Author count & Group, plus a CRO & One, plus the models & Every deposit \\
\end{tabularx}
\end{apptable}
\vspace{-0.6cm}
\tabcap{Table 5. The ten operating axes of Figure 2 as citable text,\\
with the section of this paper where each value is established.}

The rest of the paper is organized as its evidence rather than as an argument.
\S\ref{sec:methods} gives the method and the two properties that make it repeatable. \S\ref{sec:ind} through
\S\ref{sec:funding} give the four deliverables the method produced, each written to a similar
length because each is load-bearing: the IND, the two trial protocols, the
legislative framework, and the funding proposals. \S\ref{sec:peerreview} gives the review regime
that cleared all four and the measurements that justify trusting it. \S\ref{sec:limitations} states
what remains exposed and what future work consists of, and \S\ref{sec:conclusions} states what
follows. A reader who wants only the numbers can read Tables 8 through 20 and
stop.

\subsection{What this paper is, and what it is not}

This paper is a disclosure of a production method and a record of what that
method produced. It is not a clinical result, a regulatory filing, or an enacted
statute, and it makes no claim in any of those directions. The IND described in
\S\ref{sec:ind} is deposited, not filed. The protocols of
\S\ref{sec:protocol} are written, not approved. The bill of
\S\ref{sec:legislation} is drafted, not introduced. The applications of
\S\ref{sec:funding} are submitted, and none has yet been awarded.

What is claimed is narrower and checkable. First, that the four deliverables
exist, are deposited under digital object identifiers, and were produced on the
timescale the tables in this paper report. Second, that they were produced by
one author directing a single model family under a single master prompt, and
reviewed during development by two independent model manufacturers. Third, that
the review regime of \S\ref{sec:peerreview} is the only one available at this
production rate, because the prior system's arithmetic does not permit it.
Fourth, that the remaining obstacle to running the trial is institutional and
financial rather than documentary. Each of the four can be falsified by a reader
with the repository and no special access, and \S\ref{sec:ind} names the four
falsification routes for the dossier specifically.

The tone of the Federal record surveyed above is ambitious, and this paper is
written to match it rather than to hedge against it. An administration that has
pointed executive orders, appropriations and agency capability at the
intersection of artificial intelligence and cancer within eighteen months is
entitled to expect sponsors to move at a comparable rate, and the purpose of
this disclosure is to show that a sponsor can.
